\section{问题一：常热物性预热模型与结果}
\subsection{特征尺度与求解设置}
由题给物性，热扩散率 $\alpha=k/(\rho c_p)=1.689\times10^{-7}$~m$^2$/s；初始水分扩散率 $D_0=4.938\times10^{-9}$~m$^2$/s。以半径为尺度，$R_0^2/\alpha\approx2369$~s，$R_0^2/D_0\approx8.10\times10^4$~s，说明短时热扰动比水分扰动传播得快。对应 $h_TR_0/k=1.389$、$h_mR_0/D_0=3.240$，内部梯度不可忽略。
第一问保持密度、比热容与导热系数为常数，水分方程仍通过 $D(C)$ 非线性求解。每1~s积分，采样到 $0,0.1,\ldots,2.0$~cm。中心控制体为完整小圆盘，并非半圆盘。

\subsection{温度与含水率结果}
\begin{table}[!htbp]\centering
\caption{问题1温度分布（$^\circ$C）}
\begin{tabular}{lrrrrr}\toprule
$t$/s & 0 cm & 0.5 cm & 1.0 cm & 1.5 cm & 2.0 cm\\\midrule
100 & 28.0001 & 28.0004 & 28.0041 & 28.0324 & 28.1802\\
300 & 28.0414 & 28.0642 & 28.1522 & 28.3683 & 28.8495\\
600 & 28.4549 & 28.5376 & 28.8053 & 29.3168 & 30.1650\\
900 & 29.3262 & 29.4602 & 29.8774 & 30.6177 & 31.7305\\
1200 & 30.5447 & 30.7117 & 31.2241 & 32.1143 & 33.4304\\
1500 & 31.9976 & 32.1887 & 32.7681 & 33.7486 & 35.1218\\
1800 & 33.5770 & 33.7736 & 34.3654 & 35.3631 & 36.7865\\
\bottomrule\end{tabular}\end{table}

\begin{table}[!htbp]\centering
\caption{问题1含水率分布（kg/kg）}
\begin{tabular}{lrrrrr}\toprule
$t$/s & 0 cm & 0.5 cm & 1.0 cm & 1.5 cm & 2.0 cm\\\midrule
100 & 2.5500 & 2.5500 & 2.5500 & 2.5500 & 2.2475\\
300 & 2.5500 & 2.5500 & 2.5500 & 2.5492 & 2.0511\\
600 & 2.5500 & 2.5500 & 2.5500 & 2.5352 & 1.8771\\
900 & 2.5500 & 2.5500 & 2.5497 & 2.5045 & 1.7549\\
1200 & 2.5500 & 2.5500 & 2.5482 & 2.4645 & 1.6587\\
1500 & 2.5500 & 2.5499 & 2.5445 & 2.4206 & 1.5788\\
1800 & 2.5500 & 2.5497 & 2.5383 & 2.3755 & 1.5103\\
\bottomrule\end{tabular}\end{table}

1800~s时中心与表面温度分别为\QoneCenterT{}、\QoneSurfaceT{}~$^\circ$C，含水率分别为\QoneCenterC{}、\QoneSurfaceC{}~kg/kg。温度由外向内提高，水分主要从外层逸出。表层变干后 $e^{-0.89/C}$ 下降，形成低扩散率区域，解释了中心含水率在预热阶段仍接近初始值。
\begin{figure}[!htbp]\centering\includegraphics[width=0.94\textwidth]{q1_preheat_profiles.pdf}
\caption{问题一的径向温度与含水率分布}\label{fig:q1}\end{figure}
题目指定七个时刻的表格由终算数据生成，导出程序使用同一数据生成 result1.xlsx；其格式含温度、水分浓度两表，时间为1--1800~s。四位小数遵从题目格式，数值误差另列于模型检验。
